\documentclass[10pt,a4paper]{amsart}
\usepackage[margin=19mm]{geometry}
\usepackage{amsmath,amssymb,mathtools}
\usepackage[scr=rsfs,cal=euler]{mathalfa}
\usepackage{xfrac}
\usepackage[unicode,hidelinks,bookmarks=false]{hyperref}
\newcommand{\ie}{i.e.\ }
\newcommand{\cf}{cf.\ }
\newcommand{\resp}{respectively\ }
\newcommand{\Subsection}{\subsection}
\providecommand{\nobreakdash}{\nobreak\hbox{-}}
\setlength{\emergencystretch}{2em}
\multlinegap0pt
\usepackage{amssymb}
\usepackage{xcolor}
\newtheorem{proposition}{Proposition}
\title{What \texorpdfstring{$R_0$}{R0} Deletes: Eigenvectors, Non-Normality, and the Social Content of the Basic Reproduction Number}
\author{Fabio Sanchez}
\date{Source: arXiv:2607.03979v1 (CC BY 4.0)}
\begin{document}
\maketitle

\noindent\emph{Source and licence.} What \texorpdfstring{$R_0$}{R0} Deletes: Eigenvectors, Non-Normality, and the Social Content of the Basic Reproduction Number. Version arXiv:2607.03979v1 (CC BY 4.0), distributed by arXiv under the Creative Commons Attribution 4.0 International licence, http://creativecommons.org/licenses/by/4.0/, as recorded in the arXiv licence metadata for this exact version and shown on the abstract page https://arxiv.org/abs/2607.03979v1. Full licence notice: \texttt{licenses/arxiv-2607.03979-CC-BY-4.0.txt}.\par\smallskip

\noindent\textit{Editorial scope.} Counted unit: the article's proposition prop:kant (source0.tex lines 124--143). The article's complete original proof of it is delivered with it (source0.tex lines 145--151, one \texttt{proof} environment of 7 lines). No mathematics is written, repaired or re-derived: the statement and the proof are the article's own lines, in the article's own notation, and no retained line is substituted or re-flowed.  No further source-local context is needed: every cross-reference of every retained line resolves inside the two delivered spans.  Nothing was omitted from any retained span, so the package carries no omission record.  No retained line contains a citation, so the wrapper carries no bibliography and no bibliography entry.  No retained line contains a figure environment, an image-inclusion command or drawing code, so no image is delivered and the package declares no asset dependency.  Editorial formatting only, in addition: an AMS \texttt{amsart} preamble that restates the article's own declarations and macros used by the retained lines, namely the environments \texttt{proposition} on separate counters, together with the standard packages those lines need.  The retained environments print in this document as Proposition 1 in order of appearance, whereas the article prints the same items with the numbering of its own full document.  The complete native source of the article as distributed in the arXiv e-print for this exact version is bundled unchanged as \texttt{sources/204435/source0.tex}.
\begin{proposition}\label{prop:kant}
Let $C$ be symmetric, entrywise nonnegative and irreducible, and $Q=\operatorname{diag}(q_i)$ with $q_i>0$. Write $K=QC$, with Perron eigenvalue $\rho$, positive right eigenvector $v$ $(Kv=\rho v)$ and left eigenvector $w$ $(w^{\top}K=\rho w^{\top})$. Then:
\begin{enumerate}
\item[\textup{(i)}] $K$ is diagonally symmetrizable, $K=Q^{1/2}SQ^{-1/2}$ with $S=Q^{1/2}CQ^{1/2}$ symmetric; its eigenvalues are real, and $v=Q^{1/2}\varphi$, $w=Q^{-1/2}\varphi$, where $\varphi>0$ is the unit Perron eigenvector of $S$ (its dominant eigenvector, with all entries positive).
\item[\textup{(ii)}] Writing $p_i=\varphi_i^{2}$ \textup{(}a probability vector over groups\textup{)} and $\mathbb{E}_p[\cdot]$ for the corresponding weighted average, the condition number and source--burden angle satisfy
\begin{equation}
\kappa=\frac{\lVert v\rVert\,\lVert w\rVert}{w^{\top}v}
=\Big(\textstyle\sum_i p_i q_i\Big)^{1/2}\Big(\sum_i p_i q_i^{-1}\Big)^{1/2}
=\big(\mathbb{E}_p[q]\,\mathbb{E}_p[q^{-1}]\big)^{1/2},
\qquad \cos\theta=\frac1\kappa .
\end{equation}
\item[\textup{(iii)}] With $m=\min_i q_i$, $M=\max_i q_i$ and contrast $r=M/m$,
\begin{equation}
1\;\le\;\kappa\;\le\;\tfrac12\!\left(\sqrt{r}+\tfrac{1}{\sqrt{r}}\right),
\qquad\text{equivalently}\qquad
\theta\;\le\;\arccos\frac{2\sqrt{r}}{1+r}.
\end{equation}
The lower bound holds with equality iff $Q$ is a scalar matrix. The upper bound is sharp: it is attained when $q$ takes only its extreme values $\{m,M\}$ (with $C$ placing equal Perron mass on them), and is a strict supremum otherwise, since a positive Perron vector must assign positive weight to any intermediate susceptibilities.
\end{enumerate}
\end{proposition}

\begin{proof}
\textup{(i)} $Q^{1/2}SQ^{-1/2}=Q^{1/2}\bigl(Q^{1/2}CQ^{1/2}\bigr)Q^{-1/2}=QC=K$, and $S^{\top}=Q^{1/2}C^{\top}Q^{1/2}=S$. Being similar to the symmetric matrix $S$, $K$ has real eigenvalues. Since $S$ is symmetric, nonnegative and irreducible, Perron--Frobenius provides a unit eigenvector $\varphi>0$ with $S\varphi=\rho\varphi$ and $\rho=\rho(K)$. Then $K\bigl(Q^{1/2}\varphi\bigr)=Q^{1/2}S\varphi=\rho\,Q^{1/2}\varphi$, and $\bigl(Q^{-1/2}\varphi\bigr)^{\top}K=\rho\bigl(Q^{-1/2}\varphi\bigr)^{\top}$.

\textup{(ii)} $w^{\top}v=\varphi^{\top}Q^{-1/2}Q^{1/2}\varphi=\varphi^{\top}\varphi=1$; moreover $\lVert v\rVert^{2}=\varphi^{\top}Q\varphi=\sum_i q_i\varphi_i^{2}$ and $\lVert w\rVert^{2}=\varphi^{\top}Q^{-1}\varphi=\sum_i q_i^{-1}\varphi_i^{2}$. Hence $\kappa=\lVert v\rVert\,\lVert w\rVert$ takes the stated form and $\cos\theta=w^{\top}v/(\lVert v\rVert\lVert w\rVert)=1/\kappa$.

\textup{(iii)} By Cauchy--Schwarz, $\mathbb{E}_p[q]\,\mathbb{E}_p[q^{-1}]\ge\bigl(\mathbb{E}_p[\sqrt q\cdot q^{-1/2}]\bigr)^{2}=1$, with equality iff $q$ is $p$-almost-surely constant; as $\varphi>0$ this forces all $q_i$ equal. The Kantorovich inequality, $\mathbb{E}_p[q]\,\mathbb{E}_p[q^{-1}]\le (M+m)^{2}/(4Mm)$ for $q_i\in[m,M]$, gives the upper bound after writing $r=M/m$ and taking square roots; the angle form follows from $\cos\theta=1/\kappa$.
\end{proof}

\end{document}
